\section{Background and problem setup}
\label{sec:background}

\subsection{Repeated measurements and classical readout features}

We work in a reset-based or repeated-measurement QRC setting. At each labeled time step $t$, a scalar input $u_t$ is encoded into a fixed $Q$-qubit circuit, the reservoir dynamics are applied, and the circuit is measured repeatedly. In the synthetic suite the measured observables are $Z_1,\dots,Z_Q$ and $X_1,\dots,X_Q$, implemented through two measurement settings. Each repeated execution returns a shot-level measurement vector
\begin{equation}
m_t^{(n)} \in \{-1,+1\}^{M}, \qquad n=1,\dots,\Nshots,
\end{equation}
with $M=2Q$ scalar outcomes collected across $\Rsettings=2$ settings.

Temporal memory is introduced classically through a leaky integrator,
\begin{equation}
r_t^{(n)} = (1-\eta)r_{t-1}^{(n)} + \eta m_t^{(n)}, \qquad r_0^{(n)}=0,
\label{eq:leaky}
\end{equation}
followed by a tapped-delay embedding of length $L$. The resulting shot-level readout feature is
\begin{equation}
\phi_t^{(n)} = [r_t^{(n)},r_{t-1}^{(n)},\dots,r_{t-L+1}^{(n)}] \in \mathbb{R}^{d},
\end{equation}
with feature dimension $d = ML = 2QL$.

\subsection{Exact budget accounting}

The central resource in this paper is the total number of circuit executions. If each labeled time step uses $\Rsettings$ measurement settings and $\Nshots$ repeated shots per setting, then the budget is
\begin{equation}
\Stotal = \Rsettings \Nshots T,
\label{eq:budget}
\end{equation}
where $T$ is the number of retained labeled time steps before the chronological split. In the default synthetic suite $\Rsettings=2$, because one $Z$-basis and one $X$-basis experiment are executed at every time step. The factor $\Rsettings$ is therefore part of every budget-fair comparison.

After discarding a washout of $t_0$ steps, we use a chronological split. Let $\Ttrain$ denote the number of labeled time steps in the final training block before any grouping is performed. We define the EV row-to-dimension ratio as
\begin{equation}
\rhoEV = \frac{\Ttrain}{d}.
\label{eq:rhoev}
\end{equation}
Under the accounting in Eq.~\eqref{eq:budget},
\begin{equation}
\rhoEV \approx \frac{\alpha}{d}\left(\frac{\Stotal}{\Rsettings \Nshots}-t_0\right),
\label{eq:rhoevbudget}
\end{equation}
up to integer effects of the split. Throughout the paper, $\rhoEV$ is the main diagnostic for whether the EV readout is row-limited.

\subsection{Protocol endpoints and evaluation}

The three main protocol endpoints are:
\begin{enumerate}[leftmargin=*,label=(\alph*)]
    \item \textbf{EV}: average all shots at a time step into one feature vector,
    \item \textbf{Raw}: keep every shot as its own training row,
    \item \textbf{\Split{}}: partition the shots into disjoint groups and average within each group.
\end{enumerate}

All protocols ultimately return one prediction per labeled time step. Grouped predictions are therefore averaged back to the time-step level before evaluation. Our primary metric is test $\NRMSE$; we also report the standard-deviation ratio (SR) for selected plots because it quantifies the amplitude attenuation that is expected under noisy-covariate training.
